\chapter{README}

章タイトルのデザインは quotchap パッケージ\footnote{\url{https://ctan.org/pkg/quotchap}}のパクリ。

\section{執筆のしかた}

このリポジトリを \verb`git clone` したら、\verb`*.tex` をてきとーに編集して、\verb`make` するだけ。

必要に応じて、\verb`git remote remove origin` してから自前のリポジトリに \verb`git remote add origin ...` しなおすべし。

\subsection{\LaTeX 環境}

Lua\LaTeX + jlreq が必要。\TeX Live で動くはず。

docker で構築するのが楽ちん。

\begin{lstlisting}[style=tty]
% docker pull texlive/texlive:latest
\end{lstlisting}

docker を使わない場合は \verb`Makefile` の \verb`LATEXMK=...` の行を修正する必要がある。

\begin{lstlisting}
# docker を使う場合
LATEXMK= docker run --rm -v .:/workdir texlive/texlive:latest latexmk
# 使わない場合
#LATEXMK= latexmk
# ローカルではdockerを使わないが、github actions でコンパイルするときには使う
#LATEXMK= $(shell command -v latexmk || echo "docker run --rm -v .:/workdir texlive/texlive:latest latexmk")
\end{lstlisting}

\subsection{コンパイル}

\verb`make` 一発。

\verb`Makefile` をとくに修正していない場合、電子版(画面上での閲覧を想定したPDF)の \verb`main.pdf` を生成する。また、\verb`make all` ではそれに加えて印刷版(印刷所への入稿を想定したPDF)の \verb`print.pdf` を生成する。コンパイルのたびに両方生成するのは時間のムダなので、デフォルトではあえて片方だけ生成するようにしている。

\verb`make` の際、\verb`git` のコミットハッシュが記載された \verb`revid.tex` というファイルを自動で生成する。この値は \verb`\revid` というマクロで参照できる。本ドキュメントでは奥付で使用している。また、PDFのメタ情報にも埋め込んでいる。

\begin{lstlisting}[style=tty]
% git rev-parse --short main
1eb445e
% cat revid.tex
\def\revid{1eb445e}
% pdfinfo -meta main.pdf | fgrep VersionID
                  <pdfaProperty:name>VersionID</pdfaProperty:name>
      <xmpMM:VersionID>1eb445e</xmpMM:VersionID>
\end{lstlisting}

画像を読み込むための \verb`\includegraphics{}` とソースコードを読み込むための \verb`\lstinputlisting{}` で指定されているファイルは自動的に依存関係に追加されるので(\verb`Makefile` の \verb`INCLUDE=...` の行)、これらのファイルを修正すれば \verb`*.tex` を修正していなくても \verb`make` で再生成される。ただし、環境変数 \verb`$TEXINPUTS` で指定されたディレクトリまでは見ない。\verb`$TEXINPUTS` に頼らない書き方をするのがよい。

\verb`make` が使えない環境(Windows や Overleaf のような Web 上のもの)は知らん。がんばって何とかしてくれ。

\subsection{github workflow}

\verb`main` または \verb`master` ブランチで github に push すると、github 上で \verb`make all` される。docker を使わずにコンパイルするように \verb`Makefile` の \verb`LATEXMK=...` の行を修正していた場合、github でのビルドがコケるので注意。

自動コンパイルしたくない場合は \verb`.github/` 以下をまるごと削除すべし。

\subsection{印刷版と電子版}

前述のように \verb`make all` で2種類のPDFを生成する。主な違いは以下。

\begin{itemize}
\item 印刷版は見開き左右でページデザインが異なるが(twoside)、左右の概念がない電子版はどのページも同じデザイン(oneside)
\item 電子版は1ページ目に表紙画像を挿入するが、印刷版では別ファイルでの入稿を想定して挿入しない
\item 印刷版はグレースケール\footnote{ただし、カラー画像がグレースケールに変換されるわけではない}、電子版はRGBカラー
\item 印刷版ではトンボと隠しノンブルを出力する
\end{itemize}

これ以外に電子版、印刷版で異なる処理をさせたい場合は、\verb`main.tex` で以下のように記述する。

\begin{lstlisting}
\ifdefined\PRINT
  % 印刷版の処理
\else
  % 電子版の処理
\fi
\end{lstlisting}

これは\ifdefined\PRINT 印刷版\else 電子版\fi のPDFである。

\verb`print.tex` は電子版か印刷版かを区別するための \verb`\PRINT` マクロを定義してから \verb`main.tex` を読み込むだけのファイルで、通常はとくに修正する必要はない。

\section{セクション}

ここからは見た目を確認するためのデモ。

\subsection{サブセクション}

びよーん。

\subsubsection{サブサブセクション}

ぱおーん。

\begin{column}{コラムタイトル}

コラムはこう書く\footnote{コラム内で脚注を書くとこうなる}。

目次にも自動的に掲載される。

\end{column}

ソースコードはリスト\ref{src}のように書く。

\begin{lstlisting}[caption=hello.sh, label=src]
echo Hello, world!
\end{lstlisting}

枠のないバージョン。

\begin{lstlisting}[style=tty]
$ sh hello.sh
Hello, world!
\end{lstlisting}

なんか引用文を入れるのによさそうな環境。えらい人曰く、

\begin{leftbar}
腹へった。
\end{leftbar}

\subsection{フォント}

\def\sampletext{\noindent%
A quick brown fox jumps over the lazy dog. \\
色は匂へど散りぬるを我が世誰ぞ常ならむ有為の奥山今日越えて浅き夢見し酔ひもせず \\
0123456789
}

\textrm{\textbackslash{}rmfamily}
\begin{snugshade}
\rmfamily\sampletext
\end{snugshade}

\textsf{\textbackslash{}sffamily}
\begin{snugshade}
\sffamily\sampletext
\end{snugshade}

\texttt{\textbackslash{}ttfamily}
\begin{snugshade}
\ttfamily\sampletext
\end{snugshade}

